\documentclass[12pt]{article}

\usepackage[a4paper,margin=1in]{geometry}
\usepackage{amsmath,amssymb}
\usepackage{bm}
\usepackage{enumitem}
\usepackage{hyperref}
\usepackage{setspace}

\setstretch{1.15}

\title{\bf BIAS-CONDITIONED LOW-PRECISION NEURAL NETWORK TRAINING VIA RANK-ONE MEAN DECOMPOSITION}
\author{}
\date{}

\begin{document}

\maketitle

\section*{ABSTRACT}

A computer-implemented method and system for stabilizing low-precision neural network training are disclosed. An activation tensor associated with a neural network layer is decomposed into a rank-one mean component and a residual component. The rank-one mean component and residual component are independently quantized and recombined during matrix multiplication operations. In certain embodiments, a symmetric decomposition is applied to gradient tensors during backpropagation. The disclosed approach reduces dynamic range inflation and second-moment instability without requiring spectral decomposition operations such as singular value decomposition. The method is hardware-friendly and suitable for GPU, ASIC, and accelerator-based training environments.

\newpage

\section{FIELD OF THE INVENTION}

The present invention relates to machine learning systems. More specifically, the invention relates to low-precision neural network training and quantization techniques that improve numerical stability and hardware efficiency.

\section{BACKGROUND}

Neural networks, including transformer-based language models, are increasingly trained using reduced numerical precision formats such as FP8, FP4, and other low-bit floating point or blockwise quantization schemes. Reduced precision improves computational efficiency, memory utilization, and hardware throughput.

However, reduced precision training suffers from instability caused by extreme activation magnitudes and inflated adaptive second-moment statistics. Existing mitigation techniques often rely on spectral decompositions, whitening procedures, or orthogonal transformations that require significant computational overhead and are poorly aligned with high-throughput accelerator hardware.

Accordingly, there exists a need for an efficient and hardware-compatible method for mitigating bias-driven outliers and second-moment inflation during low-precision neural network training.

\section{SUMMARY OF THE INVENTION}

The invention provides a method for bias-conditioned training of neural networks. The method includes:

\begin{enumerate}
    \item Computing a feature-wise mean vector from an activation tensor.
    \item Constructing a rank-one mean tensor from said mean vector.
    \item Computing a residual activation tensor by subtracting the rank-one mean tensor.
    \item Independently quantizing the rank-one mean tensor and the residual tensor.
    \item Performing matrix multiplication using the independently quantized components.
    \item Optionally performing symmetric decomposition on gradient tensors during backpropagation.
\end{enumerate}

The disclosed approach avoids singular value decomposition and other expensive spectral operations while stabilizing low-bit training.

\newpage

\section{DETAILED DESCRIPTION}

\subsection{Activation Mean Decomposition}

Let an activation tensor be reshaped into a matrix:

\[
X \in \mathbb{R}^{l \times m}
\]

where $l$ represents batch size multiplied by sequence length, and $m$ represents feature dimension.

A feature-wise mean vector is computed as:

\[
\mu_X = \frac{1}{l} \bm{1}^\top X
\]

where $\bm{1}$ is an all-ones vector of dimension $l$.

A rank-one mean tensor is constructed:

\[
M_X = \bm{1} \mu_X^\top
\]

The residual activation tensor is:

\[
X_R = X - M_X
\]

\subsection{Independent Quantization}

Let $Q_b(\cdot)$ represent a quantization operator to a reduced precision format.

\[
\overline{M}_X = Q_b(M_X)
\]
\[
\overline{X}_R = Q_b(X_R)
\]

Independent scaling parameters may be used for $\overline{M}_X$ and $\overline{X}_R$.

\subsection{Matrix Multiplication Reconstruction}

Given a quantized weight matrix $\overline{W}$:

\[
\widehat{Y} = (\overline{M}_X + \overline{X}_R) \overline{W}
\]

In certain embodiments, the mean tensor is not explicitly materialized, and subtraction is performed on-the-fly.

\subsection{Backward-Pass Decomposition}

Given gradient tensor:

\[
D = \frac{\partial L}{\partial Y}
\]

A mean vector is computed:

\[
\mu_D = \frac{1}{l} \bm{1}^\top D
\]

Residual gradient:

\[
D_R = D - \bm{1}\mu_D^\top
\]

Independent quantization is applied:

\[
\overline{M}_D = Q_b(\bm{1}\mu_D^\top)
\]
\[
\overline{D}_R = Q_b(D_R)
\]

Gradients are computed:

\[
\widehat{\frac{\partial L}{\partial X}} = (\overline{M}_D + \overline{D}_R) \overline{W}^\top
\]

\[
\widehat{\frac{\partial L}{\partial W}} = (\overline{M}_X + \overline{X}_R)^\top (\overline{M}_D + \overline{D}_R)
\]

\subsection{Optional Embodiments}

Embodiments may include:

\begin{itemize}
    \item Partial mean scaling using a factor $\alpha$.
    \item Per-head or per-group mean computation.
    \item Mixed-precision encoding of mean and residual.
    \item Application during inference-only quantization.
    \item Distributed training implementations.
    \item Optional residual spike conditioning without full spectral decomposition.
\end{itemize}

\subsection{Hardware Implementation}

The method may be implemented using:

\begin{itemize}
    \item Reduction kernels for mean computation.
    \item Elementwise subtraction operations.
    \item Existing quantized GeMM hardware.
\end{itemize}

No singular value decomposition or eigen decomposition is required.

\newpage

\section*{CLAIM APPENDIX (DRAFT CLAIMS)}

\textbf{Claim 1.} A computer-implemented method for training a neural network using low-precision numerical representation, comprising:

\begin{enumerate}
    \item receiving an activation tensor associated with a neural network layer;
    \item computing a feature-wise mean vector from the activation tensor;
    \item constructing a rank-one mean tensor from the feature-wise mean vector;
    \item generating a residual activation tensor by subtracting the rank-one mean tensor from the activation tensor;
    \item independently quantizing the rank-one mean tensor and the residual activation tensor;
    \item performing a matrix multiplication using the independently quantized tensors and a quantized weight tensor.
\end{enumerate}

\textbf{Claim 2.} The method of claim 1, further comprising applying a corresponding decomposition and independent quantization to gradient tensors during backpropagation.

\textbf{Claim 3.} The method of claim 1, wherein the reduced numerical precision format comprises a 4-bit floating point representation.

\textbf{Claim 4.} The method of claim 1, wherein singular value decomposition is not performed.

\textbf{Claim 5.} A neural network training system configured to perform the method of claim 1.

\textbf{Claim 6.} A non-transitory computer-readable medium storing instructions that cause one or more processors to perform the method of claim 1.

\end{document}